\documentclass[instructions.tex]{subfiles}
\begin{document}
 
\chapter{On peut facilement pratiquer la pénitence dans tous les états.}
 
Presque tous les hommes vivent dans l’exercice de la pénitence; mais parce qu’ils n’en font pas un saint usage, la plupart meurent sans l’avoir faite. On croit la pénitence bien éloignée de nous, tandis que nous en sommes environnés. Y a-t-il une condition qui n’ait ses peines, ses embarras, ses tentations, ses épreuves? Qu’on est aveugle de ne pas profiter de la bonté de Dieu, qui accepte comme une pénitence méritoire ce qui nous arrive même par nécessité!
 
Que vous coûterait-il, dans vos afflictions, dans vos délaissements, de dire avec Jésus-Christ : Mon père, que votre volonté soit faite? Que vous coûterait-il, lorsque vous perdez votre santé, ou qu’un injuste usurpateur vous enlève vos biens, de dire avec Job : Dieu me l’avait donné, Dieu me l’a ôté, que son saint nom soit béni? Si nous avons reçu des faveurs de sa main libérale, n’est-il pas juste que nous recevions aussi les maux de sa main paternelle\footnote{Si bona suscepimus de manu Domini, mala quare non suscipiamus? Job,} ? En souffrant pour Dieu, en souffrirait-on davantage?
 
Combien de choses dont la privation vous est facile, et qui vous seraient d’un grand mérite! Seriez-vous fort incommodé de vous priver d’une lecture profane, pour lire un livre de piété; de fréquenter une compagnie sainte, au lieu de fréquenter des gens dissipés et mondains; d’aller visiter Jésus-Christ, au lieu d’aller au spectacle? Seriez-vous fort incommodé de vous habiller avec plus de modestie, de parler avec plus de retenue, de vous modérer dans vos repas, de donner aux pauvres ce que vous donnez pour le jeu ou pour le luxe? Oh! que de mérites perdus par votre faute!
 
La Providence nous fournit, à toute occasion, mille moyens de faire pénitence : un petit soulagement qu’on refuse à son corps, une situation gênante, une incommodité passagère qu’on souffre sans se plaindre; l’humeur dure et fâcheuse d’une personne avec laquelle nous sommes obligés de vivre, que nous supportons avec charité; une parole désagréable, un reproche humiliant dont on fait le sacrifice en silence; un peu de froid ou de chaud qu’on souffre avec patience; un peu de faim ou de soif qu’on ne soulage pas d’abord; un service qu’on rend au prochain; une délicatesse, un morceau, un fruit, un assaisonnement dont on se prive, et autres semblables mortifications, ne sont pas indignes d’être offerts à Dieu.
 
Tout cela paraît peu de chose; mais c’est quelque chose de grand, dit saint Augustin, d’être fidèle à Dieu dans les plus petites choses\footnote{In minimis autem fidelis esse magnum est. S. Aug.} . Ce ne sont là que comme de petites gouttes, mais qui peuvent peu à peu remplir le calice amer que Jésus-Christ nous a laissé à boire. Par ces petites pratiques, que les mondains regardent comme des riens, nous pouvons payer beaucoup, et grossir notre trésor pour l’éternité, à l’exemple de la fourmi, qui fait ses provisions de ce qu’on foule aux pieds. En offrant à Dieu ces fréquens et petits sacrifices, ayons de la confusion de lui offrir si peu de chose; mais adorons sa bonté qui s’en contente. C’est en méprisant ces petites choses, que tant de gens vivent sans pénitence, et se perdent.
 
Quel repentir à la mort et dans l’éternité, d’avoir laissé échapper tant d’occasions d’effacer ses péchés!
 
\section{Résolutions}
 
\primo J’offrirai tous les jours à Dieu, en satisfaction pour mes péchés, les peines de mon état.
 
\secundo J’accepterai de même toutes les afflictions qui pourront m’arriver, de quelque source qu’elles me viennent. 

\prayer{Mon Dieu, j’unis toutes mes souffrances à celles qu’éprouva autrefois mon Sauveur; agrééz-les comme des satisfactions que j’offre à votre justice pour l’expiation de mes nombreux péchés.}
 
\end{document}
